
\subsection{扰动后的模空间逸出仓西邻域 II}
\label{subsec:7-2-8}

本小节利用第~\ref{subsec:7-2-3} 小节的思路，
完成第 3 章所述命题~\ref{prop:3-5-2}
和定理~\ref{thm:3-5-11} 的证明。
定理~\ref{thm:7-2-72}
与命题~\ref{prop:7-2-78}
将起关键作用。

先在障碍丛上固定适当的范数，
用以度量两个多值截面
$\mathfrak s_{\mathfrak d,k,\beta,\vec P}$
之间的距离。
还要固定这些模空间的仓西邻域：
仓西结构在 $X$ 上的定义，
先验地只为每个 $p\in X$ 指定一个仓西邻域的芽；
这里为这些芽选定代表。
并假定所选邻域形成\emph{良好坐标系}，
其意义见 \textbf{[FuOn99II]} 定义 6.1
或本书附录 A1 的引理 A1.11。
这样就能用该坐标系归纳构造所需的多值截面。

取两族以 $\lambda$ 为参数的正数
$\{\mathcal K_\lambda\}_\lambda$
和 $\{\delta_\lambda\}_\lambda$，使
\[
\mathcal K_\lambda\nearrow\infty,
\qquad
\delta_\lambda\searrow0
\quad(\lambda\to\infty).
\]
由命题~\ref{prop:7-2-35}，
对每个 $\lambda$，有可数集
$\mathcal X_{(g)}(L,\lambda)$
（$g\le\mathcal K_\lambda$），
它生成 $C_{(g)}(L,\lambda)$。
为与第 3.5 节的记号一致，
或许把 $\mathcal X_{(g)}(L,\lambda)$
写成 $\mathcal X_{1,(g)}(L,\lambda)$ 更合适；
这里为简便仍用前者。
当 $\|(\mathfrak d,\beta)\|\le\mathcal K_\lambda$ 时，
还存在多值截面族
$\mathfrak s_{\mathfrak d,k,\beta,\vec P,\lambda}$，
它以 $\mathcal K=\mathcal K_\lambda$、
$\delta=\delta_\lambda$ 满足 (7.2.36)。

命题~\ref{prop:7-2-78} 表明，
上述多值截面族可选得使
\[
g_0(\lambda),g_1(\lambda)\nearrow\infty,
\qquad K(\lambda)\nearrow\infty,
\qquad n(\lambda)\nearrow\infty
\quad(\lambda\to\infty),
\]
并满足以下条件：

\noindent\textup{(7.2.91.1)}\quad
$C_{(g_1(\lambda))}(L,\lambda)$ 上
存在滤过 $A_{n(\lambda),K(\lambda)}$ 结构
$\mathfrak m^\lambda_{k,\beta}$。

\noindent\textup{(7.2.91.2)}\quad
在 $C_{(g_0(\lambda))}(L,\lambda)$ 上，
$\mathfrak m^\lambda_{k,\beta}
=\mathfrak m^{\lambda,\mathrm{geo}}_{k,\beta}$；
这里的 $\mathfrak m^{\lambda,\mathrm{geo}}_{k,\beta}$
如式~\eqref{eq:7-2-77}。

此后简记
\begin{align}
C(L,\lambda)
&=C_{(g_1(\lambda))}(L,\lambda),
\tag{7.2.92.1}\label{eq:7-2-92-1}\\
C_0(L,\lambda)
&=C_{(g_0(\lambda))}(L,\lambda).
\tag{7.2.92.2}\label{eq:7-2-92-2}
\end{align}

\begin{proposition}{7.2.93}\label{prop:7-2-93}
在上述情形下，存在一列指标对
$(n'(\lambda),K'(\lambda))$，使
\[
(n'(\lambda),K'(\lambda))
\lesssim(n(\lambda),K(\lambda)),
\qquad
\lim_{\lambda\to\infty}
(n'(\lambda),K'(\lambda))=(\infty,\infty),
\]
并存在依赖于 $\lambda$ 的常数 $c(\lambda)$，
使当 $\lambda'>c(\lambda)$ 时，
滤过 $A_{n(\lambda),K(\lambda)}$ 代数
$(C(L,\lambda),\mathfrak m^\lambda_k)$
与滤过 $A_{n(\lambda'),K(\lambda')}$ 代数
$(C(L,\lambda'),\mathfrak m^{\lambda'}_k)$
为滤过 $A_{n'(\lambda),K'(\lambda)}$ 同伦等价。
\end{proposition}

命题~\ref{prop:7-2-93}
由定理 4.6.1 与 4.6.25 的
$A_{n,K}$ 版本推出；
它们将在后面两个小节证明。
命题~\ref{prop:7-2-93} 本身的证明
放在第 7.2.11 小节。

以下假定命题~\ref{prop:7-2-93}，
并结合定理~\ref{thm:7-2-72}
与引理~\ref{lem:7-2-69}，
完成滤过 $A_\infty$ 代数的构造。
取一列参数
\[
\lambda_0<\lambda_1<\cdots\longrightarrow\infty,
\qquad
\lambda_{i+1}>c(\lambda_i),
\]
其归纳选择如命题~\ref{prop:7-2-93} 所述。
我们将对 $i$ 归纳，
在 $C(L,\lambda_0)$ 上构造
滤过 $A_{n'(\lambda_i),K'(\lambda_i)}$ 结构
$\mathfrak m^{(\lambda_i)}$，
并对每个 $\lambda_i\ge\lambda_0$
构造滤过 $A_{n'(\lambda_i),K'(\lambda_i)}$
同伦等价
\[
\mathfrak h^{\lambda_0,\lambda_i}\colon
C(L,\lambda_0)\longrightarrow C(L,\lambda_i).
\]

作为在 $\lambda_i$ 处的归纳假设，
设已构造上述
$\mathfrak m^{(\lambda_i)}$
和 $\mathfrak h^{\lambda_0,\lambda_i}$；
并且当 $j<i$ 时，
$C(L,\lambda_0)$ 上的滤过
$A_{n'(\lambda_j),K'(\lambda_j)}$ 结构
$\mathfrak m^{(\lambda_j)}$
等于滤过 $A_{n'(\lambda_i),K'(\lambda_i)}$ 结构
$\mathfrak m^{(\lambda_i)}$ 的限制。

命题~\ref{prop:7-2-93} 给出滤过
$A_{n'(\lambda_i),K'(\lambda_i)}$ 同伦等价
\[
\mathfrak h^{\lambda_i,\lambda_{i+1}}\colon
C(L,\lambda_i)\longrightarrow C(L,\lambda_{i+1}).
\]
因 $C(L,\lambda_{i+1})$ 有滤过
$A_{n'(\lambda_{i+1}),K'(\lambda_{i+1})}$ 结构，
定理~\ref{thm:7-2-72} 表明，
可把 $C(L,\lambda_i)$ 上的滤过
$A_{n'(\lambda_i),K'(\lambda_i)}$ 结构
扩张为滤过
$A_{n'(\lambda_{i+1}),K'(\lambda_{i+1})}$ 结构，
同时把映射
$\mathfrak h^{\lambda_i,\lambda_{i+1}}$
扩张为这一较高阶的同伦等价。
再利用归纳假设中的
$\mathfrak h^{\lambda_0,\lambda_i}$，
把 $\mathfrak m^{(\lambda_i)}$
扩张为滤过
$A_{n'(\lambda_{i+1}),K'(\lambda_{i+1})}$ 结构
$\mathfrak m^{(\lambda_{i+1})}$，
并将 $\mathfrak h^{\lambda_0,\lambda_i}$
扩张为同阶的同伦等价。
定义
\[
\mathfrak h^{\lambda_0,\lambda_{i+1}}
=\mathfrak h^{\lambda_i,\lambda_{i+1}}
\circ\mathfrak h^{\lambda_0,\lambda_i}.
\]

这样对每个 $i$ 都得到所需的
$\mathfrak m^{(\lambda_i)}$
和 $\mathfrak h^{\lambda_0,\lambda_i}$。
由引理~\ref{lem:7-2-69}，
$C(L,\lambda_0)$ 上遂有滤过 $A_\infty$ 结构。
这给出第 3.5 节所需的可数生成子模
$C(L;\mathbb Q)$。
因此，在命题~\ref{prop:7-2-93} 成立的前提下，
我们已在
\[
C(L;\Lambda_{0,\mathrm{nov}})
=C(L;\mathbb Q)\widehat\otimes
\Lambda_{0,\mathrm{nov}}
\]
上构造滤过 $A_\infty$ 代数结构。

\begin{remark}{7.2.94}\label{rem:7-2-94}
上面的滤过 $A_\infty$ 代数构造
从 $\mathcal X_0(L)$ 出发，
一次性构造整个滤过 $A_\infty$ 代数。
但在定理~\ref{thm:3-1-5} 中，%
\translatornote{原书此处误引“定理 3.1.10”；
本书 (3.1.10) 是 Maurer--Cartan 方程而非定理。
所述从无滤过结构构造滤过代数的目标见定理 3.1.5，
其具体构造见定理 3.5.11；参见疑误 A-0442。}
我们先假定存在定理~\ref{thm:3-4-8}
那样的无滤过 $A_\infty$ 代数，
再主张纳入伪全纯圆盘的贡献，
以得到滤过 $A_\infty$ 代数。
为证明这一说法，
可如下修改命题~\ref{prop:7-2-35}。
假定已有 $\mathcal X_{00}(L)$，
以及在
\[
\mathcal M^{\mathrm{main},\mathfrak d_0}_{k+1}
(\beta_0;P_1,\ldots,P_k),
\qquad P_i\in\mathcal X_{00}(L),
\]
上的一个多值截面系，
它定义一个无滤过 $A_\infty$ 代数。
在此附加假设下，
仍可构造满足命题~\ref{prop:7-2-35}
全部性质的 $\mathcal X_g(L)$，
并使 $\mathcal X_0(L)\supset\mathcal X_{00}(L)$。
而且，当 $P_i\in\mathcal X_{00}(L)$ 时，
上述模空间上的多值截面
与预先给定者一致。
这些结论可以用相同方法证明，
继而以相同方式证明定理~\ref{thm:3-1-5}。
\end{remark}
